% figure_sheet_two.tex - every figure of a TWO-column paper at print size (elsarticle 5p).
%
% 1. Draw the figures with the default layout:   python fig_examples.py
%    (they are written to ../images/paper/)
% 2. Compile twice (the second run resolves the Fig. references):
%        cd latex
%        pdflatex figure_sheet_two
%        pdflatex figure_sheet_two
%
% The figures are drawn for one column: \columnwidth = 252 pt = 3.5 in (docs/paper_figure_guide.md,
% section 2). Add your own figures below, in the same way.
\documentclass[5p,twocolumn]{elsarticle}
\usepackage{graphicx,placeins}
\graphicspath{{../images/paper/}}

\begin{document}
\typeout{WIDTHS: textwidth=\the\textwidth, columnwidth=\the\columnwidth, linewidth=\the\linewidth}

\begin{frontmatter}
\title{Figure sheet, two columns}
\author{}
\begin{abstract}
Every figure of the paper at print size, one column wide.
\end{abstract}
\end{frontmatter}

% signal and residual
\begin{figure}[t]\centering
  \includegraphics[width=\columnwidth]{fig_example_snapshot.pdf}
  \caption{Signal and residual at one instant.}\label{fig:snapshot}
\end{figure}

\begin{figure}[t]\centering
  \includegraphics[width=\columnwidth]{fig_example_bands_ramp.pdf}
  \caption{A metric with its distribution as a transparency ramp.}\label{fig:ramp}
\end{figure}

\begin{figure}[t]\centering
  \includegraphics[width=\columnwidth]{fig_example_bands_patch.pdf}
  \caption{The same as a patch: 25--75\,\%.}\label{fig:patch}
\end{figure}

\begin{figure}[t]\centering
  \includegraphics[width=\columnwidth]{fig_example_tiles.pdf}
  \caption{One tile per model.}\label{fig:tiles}
\end{figure}

\begin{figure}[t]\centering
  \includegraphics[width=\columnwidth]{fig_example_split.pdf}
  \caption{Two windows of a signal.}\label{fig:split}
\end{figure}

\begin{figure}[t]\centering
  \includegraphics[width=\columnwidth]{fig_example_convergence.pdf}
  \caption{Curves with a zoom.}\label{fig:conv}
\end{figure}

\begin{figure}[t]\centering
  \includegraphics[width=\columnwidth]{fig_example_map.pdf}
  \caption{A path coloured by a value.}\label{fig:map}
\end{figure}

\FloatBarrier
\end{document}
